\section{A linear gamma outer model and the regulator layer}\label{sec:gamma-outer}
The exact gamma completion gives a static approximation on compact
positive-time intervals and on a joint time/height range.  Its finite
phase formula also resolves the regulator transition near $t=64$.
The estimates in this section are those of Notes~27--28
\cite{n27,n28}; they use the unchanged physical regulator of
Section~\ref{sec:protocol} and fixed spatial averaging, without a time
average.  The stronger spatial coverage in Section~\ref{sec:full-saddle}
will use a different gamma model.

Fix $\ell,h>0$, $|y|\le Y$, and the kernel \eqref{eq:kernel}.
For $X>0$ use the full holomorphic buffer
\[
 \Omega_X=\{z:|\Re z-X|\le\ell+2,\quad |\Im z|\le Y+h+2\}.
\]
Enlarge its fixed margin once if necessary so that every collar meeting
an averaging support lies wholly inside the buffer.  Constants may
depend on these fixed widths, $Y$ and the regulator parameter $\eta$;
they are independent of $t,X$ on the declared joint range.  The compact
statements have their own interval-dependent constants.  All onsets below
are eventual existence statements, without a numerical certificate.

\subsection{The exact convolution and a two-carrier field budget}
For every $t>0$,
\begin{equation}
 H(t,z)=\frac1{\sqrt{\pi t}}\int_{\R}e^{-v^2/t}
                    \frac{\xi(s+v)}8\,dv,
 \qquad s=\frac{1+iz}{2}.
 \label{eq:xi-heat-convolution}
\end{equation}
Indeed the Gaussian average of $\cos((z-2iv)u)$ is
$e^{tu^2}\cos(zu)$, with variance $t/2$ in the $v$ coordinate.
The complete source is integrable against
$e^{t_bu^2+(Y+h+2)u}$ on any compact time interval, justifying the
interchange.  This is the all-positive-time identity in
\cite[Section~4, equation~(35)]{polymath2019heat}; its separate high-zero
theorem for $0<t\le1/2$ is not invoked at the times considered here.

Define the exact gamma factor and its logarithmic derivative by
\begin{align}
 \mathcal M(w)&=\frac{w(w-1)}{16}\pi^{-w/2}\Gamma(w/2),&
 \alpha(w)&=\frac{\mathcal M'(w)}{\mathcal M(w)}
 =\frac1w+\frac1{w-1}-\frac{\log\pi}{2}+\frac12\psi(w/2),\notag\\
 \mathcal G_t(w)&=\mathcal M(w)e^{t\alpha(w)^2/4},&
 P_{\rm lin}(t,z)&=\mathcal G_t(s)+\mathcal G_t(1-s).
 \label{eq:gamma-linear}
\end{align}
Here $\psi=\Gamma'/\Gamma$.  The factor $1/16$ is fixed by
$\xi(w)/8=\mathcal M(w)\zeta(w)$, and $\mathcal M$ is distinct from
the real source mass $M_0$.  The inner drift $\mu'/\mu$ is local to
Section~\ref{sec:explicit}; $\alpha$ here always means the function
in \eqref{eq:gamma-linear}.  The separate gamma factors are
holomorphic and nonzero on the far-height buffers where used; no entire
extension of this local model is asserted.  It is useful to record
\begin{equation}
 D_t(w)=\alpha(w)+\frac t2\alpha(w)\alpha'(w),\qquad
 (P_{\rm lin})_z=\frac i2\{
 \mathcal G_t(s)D_t(s)-\mathcal G_t(1-s)D_t(1-s)\}.
 \label{eq:linear-derivative}
\end{equation}

\begin{proposition}[Carrier-relative linear field error]
\label{prop:linear-field-budget}
On $\Omega_X$, for $32\le t\le\tfrac14\log X$ and sufficiently
large $X$,
\begin{equation}
 |H-P_{\rm lin}|\le C\mathcal E(t,X)
       \{ |\mathcal G_t(s)|+|\mathcal G_t(1-s)|\},
 \label{eq:linear-field-error}
\end{equation}
where, for a fixed finite $C_Y$,
\begin{align}
 \mathcal E(t,X)={}&\frac{e^t(t+t^2\log^2X)}X
       +X^{-t\log2/8}+X^{C_Y}e^{-t\log^2X/128}
       +e^{t-X^{2/3}/(16t)}\notag\\
 &\le C X^{-3/4}(1+\log^4X).
 \label{eq:linear-four-payments}
\end{align}
For each fixed finite $T\ge64$ the same carrier-relative error on
$32\le t\le T$ is at most $C_T(1+\log^2X)/X$.
In particular, on every compact $I=[t_a,t_b]\Subset(64,\infty)$,
\begin{equation}
 |H-P_{\rm lin}|\le C_I\delta(X)
       \{ |\mathcal G_t(s)|+|\mathcal G_t(1-s)|\},\qquad
 \delta(X)=\frac{1+\log^2X}{X}.
 \label{eq:linear-compact-field}
\end{equation}
The denominator is the sum of carrier amplitudes, not the possibly
vanishing modulus $|P_{\rm lin}|$.
\end{proposition}
\begin{proof}
Write $\sigma=\Re s$. Split \eqref{eq:xi-heat-convolution} at
$v=2-\sigma$ and $v=-1-\sigma$.  On the right use
$\xi(w)/8=\mathcal M(w)\zeta(w)$; reflect the left by
$\xi(w)=\xi(1-w)$ and $v\mapsto-v$.  This retains two right-arm
problems centered at $w=s,1-s$, with bounded lower endpoints
$2-\sigma,1+\sigma$, and the full middle strip of real parts $[-1,2]$.

Sectorial log-gamma and digamma asymptotics
\cite[5.11.1, 5.11.2, 5.11.9 and 5.11(ii)]{DLMFGammaOuter2026},
on a slightly larger sector, give
\begin{equation}
 \begin{aligned}
 a:=\Re\alpha(w)&=\tfrac12\log(X/(4\pi))+O_Y(X^{-1}),
 &b:=\Im\alpha(w),\quad |b|\le2,\\
 \alpha'&=O_Y(X^{-1}),&\alpha''&=O_Y(X^{-2}).
 \end{aligned}
 \label{eq:linear-sector-bounds}
\end{equation}
Cauchy's formula on disks of radius a fixed fraction of $X$ supplies
the derivatives of the holomorphic remainders.  No real-axis remainder
is differentiated.  With $V=X^{1/3}$ and $v_0\le v\le V$,
\[
 \mathcal M(w+v)=\mathcal M(w)e^{v\alpha(w)+E(w,v)},\qquad
 |E(w,v)|\le C_Yv^2/X,\qquad |e^{E(w,v)}-1|\le C_Yv^2/X.
\]
The complete linear Gaussian integral is $e^{t\alpha(w)^2/4}$.
Its absolute integral divided by the modulus of that answer costs
$e^{tb^2/4}\le e^t$; the real tilted Gaussian has mean $ta/2$,
variance $t/2$, and second moment $t/2+t^2a^2/4$.
Thus the local gamma expansion pays the first term of
\eqref{eq:linear-four-payments}.  On fixed $I$ the same calculation
pays $C_I\delta(X)$.

All other source pieces are retained.  Absolute convergence of the
ordinary full Dirichlet series \cite[25.2.1]{DLMFZetaOuter2026} gives
\[
 |\zeta(w+v)-1|\le3\,2^{-\Re w-v}\qquad(\Re(w+v)\ge2).
\]
Changing the Gaussian tilt from $a$ to $a-\log2$ costs the factor
\begin{equation}
 \exp\{-ta\log2/2+t((\log2)^2+b^2)/4\}
 \le X^{-t\log2/8}.
 \label{eq:linear-arithmetic-payment}
\end{equation}
The constant terms linear in $t$ are absorbed by a fixed fraction of
$t\log X$, uniformly after the onset.  On a fixed $I$ one may instead
keep the sharper bound $C_I X^{-t_a\log2/4}$, whose exponent exceeds
one.  No finite list of arithmetic modes replaces this tail.

In the bounded middle, reflect into real parts $[1/2,2]$ and use Euler
summation with $N=\lceil X\rceil$
\cite[25.2.8]{DLMFZetaOuter2026}.  It gives the sufficient bounds
\begin{equation}
 |\xi(s+v)/8|\le C_YX^4e^{-\pi X/8},\qquad
 \log|\mathcal G_t(w)|=-\pi X/8+
       \frac t{16}\log^2(X/(4\pi))+O_{I,Y}(1+\log X)
 \label{eq:linear-middle-carrier}
\end{equation}
on compact $I$.  In the joint calculation, $a\ge\log X/4$ eventually.
Completing the square pays both the bounded middle and the missing
lower linearized Gaussian tail by
$CX^{C_Y}e^{-t\log^2X/128}$.  For fixed $I$ their relative allowance
can be written
\[
 C_{I,Y}\exp\{-t_a\log^2(X/(4\pi))/16+C_{I,Y}(1+\log X)\},
\]
which is smaller than every fixed inverse power after an
interval-dependent onset.

Finally the infinite right tails are controlled on the entire ray:
\begin{equation}
 \log\frac{|\mathcal M(w+v)|}{|\mathcal M(w)|}
       \le C_Yv+\tfrac12v\log(X+v)\qquad(v\ge0).
 \label{eq:linear-ray-bound}
\end{equation}
For $v\ge V$ its right side is at most $v^2/(2t)$ on the joint range,
since $X^{1/3}/\log^2X\to\infty$ and the corresponding ray ratio is
eventually increasing.  The true and linearized tails therefore pay
$Ce^{t-X^{2/3}/(16t)}$ relative to their carriers.  On fixed $I$ they
pay $C_Ie^{-X^{2/3}/(4t_b)}$.  The carrier's height exponential has
been retained in these divisions.  Summing both reflected arms and
all four payments proves the estimates.  Since $e^t\le X^{1/4}$ and
$t\ge32$, the displayed joint rate follows; this analytical floor is
not a certified asymptotic onset.
\end{proof}

\subsection{Phase geometry and transfer for an arbitrary positive regulator}
Choose conjugate-compatible logarithms on the far-right buffer and set
\begin{align}
 \mathcal A_{\rm lin}(t,z)&=
 \exp\{(\log\mathcal G_t(s)+\log\mathcal G_t(1-s))/2\},\notag\\
 \Theta_{\rm lin}(t,z)&=
 \frac{\log\mathcal G_t(s)-\log\mathcal G_t(1-s)}{2i},\qquad
 P_{\rm lin}=2\mathcal A_{\rm lin}\cos\Theta_{\rm lin}.
 \label{eq:linear-phase}
\end{align}
The carrier is nonvanishing and positive on the real axis.  Put
\begin{align}
 A_X&=\mathcal A_{\rm lin}(t,X)
      =|\mathcal G_t((1+iX)/2)|,\qquad
 \theta=\Theta_{\rm lin}(t,X),\notag\\
 \omega&=\partial_x\theta
  =\tfrac12\Re\{\alpha(s_X)+\tfrac t2\alpha(s_X)\alpha'(s_X)\},
       \qquad s_X=(1+iX)/2.
 \label{eq:linear-frequency}
\end{align}
Uniformly on the joint range and full buffer,
\begin{equation}
 \omega=\tfrac14\log(X/(4\pi))+O(t\log X/X),\qquad
 (\Theta_{\rm lin})_{zz}=O(X^{-1}),\qquad
 |(\log\mathcal A_{\rm lin})_z|\le C.
 \label{eq:linear-phase-bounds}
\end{equation}
The last estimate uses cancellation of the leading logarithms in the
difference of the two reflected logarithmic derivatives.  Bounding
them separately by $O(\log X)$ would lose this amplitude comparison.
In particular $|\mathcal A_{\rm lin}|\asymp A_X$ on the whole buffer.
For fixed $I$ the frequency remainder is $O_I(\log X/X)$.

\begin{lemma}[Complete simple-root collars]\label{lem:linear-collars}
Let $e=C\mathcal E(t,X)$ on the joint range, or
$e=C_T(1+\log^2X)/X$ on $[32,T]$.  Fix a sufficiently small phase
radius $0<r_0<\pi/10$.  Every relevant real phase root $d_j$, defined by
$\Theta_{\rm lin}(t,d_j)=(j+1/2)\pi$, has a disjoint collar of radius
comparable to $r_0/\omega$ containing exactly one zero $a_j$ of $H$.
That zero is real and simple, and
\begin{equation}
 |a_j-d_j|\le Ce/\omega.
 \label{eq:linear-root-shift}
\end{equation}
These are all zeros that can meet the averaging buffer; at most
$C\omega$ collars meet its fixed enlargement.  For each of
$F=H,P_{\rm lin}$, its actual-root collars have the physical floors
\begin{equation}
 |F|\ge cA_X\quad\hbox{off the collars},\qquad
 |F(z)|\ge cA_X\omega|z-a_F|\quad\hbox{inside its collar}.
 \label{eq:linear-collar-floors}
\end{equation}
\end{lemma}
\begin{proof}
Divide \eqref{eq:linear-field-error} by the nonzero common carrier:
\[
 |H/\mathcal A_{\rm lin}-2\cos\Theta_{\rm lin}|
 \le Ce\{ |e^{i\Theta_{\rm lin}}|+|e^{-i\Theta_{\rm lin}}|\}.
\]
On $|z-d_j|=r_0/\omega$, phase curvature in
\eqref{eq:linear-phase-bounds} gives
$\Theta_{\rm lin}(z)-\Theta_{\rm lin}(d_j)
=\theta'(d_j)(z-d_j)+O(1/(X\omega^2))$.
The cosine is bounded below there by a fixed multiple of $r_0$.
Rouch\'e's theorem yields exactly one zero counted with multiplicity;
conjugation makes it real, and it is simple.  The identity
$|\cos(u+iv)|^2=\cos^2u+\sinh^2v$ forces any zero in the buffer to
have phase distance $O(e)$ from one of these real phase roots.
The phase derivative differs from the positive $\omega$ by $O(X^{-1})$
on the fixed rectangle.  Integrating its real part along a segment
gives an injective phase coordinate with inverse Lipschitz constant
at most $2/\omega$.  This proves completeness and
\eqref{eq:linear-root-shift}; phase spacing gives the collar count.

Recenter the collars on the actual roots, using a slightly enlarged
initial phase circle to retain boundary margins.  The trigonometric
lower bound supplies the exterior floor for both fields.  Inside a
collar, $F/(z-a_F)$ is holomorphic and zero-free, and its boundary
modulus is at least $cA_X\omega$.  The maximum principle for its
reciprocal gives the interior floor.  All disks meeting the kernel
support are full disks inside the prescribed buffer.
\end{proof}

\begin{lemma}[A collar estimate uniform over positive $Q$]
\label{lem:uniform-q-collar}
On a fixed phase-scaled disk $D$, suppose
\[
 f(u)=uB(u),\qquad g(u)=(u-e_0)C(u),\qquad |e_0|\le Ce,
\]
where the holomorphic zero-free units $B,C$ have fixed upper and
positive lower bounds, and $|B-C|\le Ce$.  For every positive function
$Q(u)$, with no derivative or positive lower bound required,
\begin{align}
 &\int_D\left|\log\frac{|g(u)|^2+Q(u)}{|f(u)|^2+Q(u)}\right|dA_u
 \notag\\[-2pt]
 &\quad+\int_D\left|\frac{Q(u)}{|g(u)|^2+Q(u)}
             -\frac{Q(u)}{|f(u)|^2+Q(u)}\right|dA_u\le Ce.
 \label{eq:uniform-q-collar}
\end{align}
In the geometry of Lemma~\ref{lem:linear-collars}, this gives on the
full averaging buffer $\Omega$, with the same physical $q$ in both fields,
\begin{equation}
 \|S_H-S_{P_{\rm lin}}\|_{L^1(\Omega)}
       +\|w_H-w_{P_{\rm lin}}\|_{L^1(\Omega)}\le Ce.
 \label{eq:linear-L1-transfer}
\end{equation}
The local estimate applies to other acquired simple-root collars with
the same unit hypotheses.  It is not a theorem about multiple events.
\end{lemma}
\begin{proof}
First replace $C$ by $B$.  Their relative error $O(e)$ changes each
integrand by $O(e)$, uniformly in $Q$.  Put $q_0(u)=Q(u)/|B(u)|^2$.
For the remaining root translation let $0\le s\le1$ vary the root
from $0$ to $e_0$, holding $q_0(u)$ fixed at each spatial point.
The absolute segment derivatives of the logarithm and rational weight
are bounded respectively by
\begin{equation}
 \frac{2|e_0|}{|u-se_0|},\qquad
 \frac{|e_0|}{2|u-se_0|}.
 \label{eq:uniform-q-derivatives}
\end{equation}
For the second use
$q_0|u-se_0|/(|u-se_0|^2+q_0)^2\le1/(4|u-se_0|)$.
The planar function $|u-se_0|^{-1}$ has a uniformly bounded integral
on the fixed disk.  Integrate first in $u$, then in $s$; the segment
itself has area zero.  This proves \eqref{eq:uniform-q-collar} without
any regularity assumption on $Q$ beyond that needed for its integrals.

For the fields take $u=\omega(z-d_j)$ and divide both by the same
nonzero $\mathcal A_{\rm lin}$.  Then
$P_{\rm lin}/\mathcal A_{\rm lin}=uB(u)$ and
$H/\mathcal A_{\rm lin}=(u-e_0)C(u)$, with $|e_0|=O(e)$.
On the collar boundary the field and root errors imply $|B-C|\le Ce$;
after filling the removable points, the maximum principle extends
that bound to the disk.  Apply the local result with
$Q=q/|\mathcal A_{\rm lin}|^2$.  Restoring area costs $\omega^{-2}$
per disk and there are $O(\omega)$ disks.  Off the collars the relative
field error is $O(e)$, so the logarithm and weight differences are
pointwise $O(e)$, independently of $q$.  This proves
\eqref{eq:linear-L1-transfer}.  The carrier logarithm cancels only in
the difference; the physical $q$ and potential offsets are unchanged.
\end{proof}

\begin{theorem}[Uniform static linear-model comparison]
\label{thm:linear-static}
For the common physical observables $R_F,U_F,W_F$ and fixed buffer above,
there exist $C,X_0$ such that
\begin{align}
 &|R_H-R_{P_{\rm lin}}|
 +\sum_{|\gamma|\le4}\bigl(
 |\partial^\gamma(U_H-U_{P_{\rm lin}})|
 +|\partial^\gamma(W_H-W_{P_{\rm lin}})|\bigr)
 \notag\\
 &\hspace{35mm}\le C X^{-3/4}(1+\log^4X)
 \quad(32\le t\le\tfrac14\log X,\quad |y|\le Y,\quad X\ge X_0).
 \label{eq:linear-static-joint}
\end{align}
For each fixed finite $T\ge64$, on $32\le t\le T$ the right side is
instead $C_T(1+\log^2X)/X$ after $X_0(T)$.  Here $\gamma$ is spatial;
the potential retains its physical additive normalization.  Both bounds
include $t=64$.  The same conclusions hold at negative heights with
$X=|x|$.
\end{theorem}
\begin{proof}
Lemma~\ref{lem:holomorphic-weak-density} needs no heat equation and
therefore applies to $P_{\rm lin}$.  Subtract its identities for the
two fields.  With $b_{r,1}=\|\nabla r\|_\infty$,
$b_{r,2}=\|\Delta r\|_\infty$, and
$|a(1-a)-b(1-b)|\le|a-b|$ on $[0,1]$, it gives
\begin{equation}
 |R_H-R_{P_{\rm lin}}|
 \le Ce(1+b_{r,1}+b_{r,2}+b_{r,1}^2).
 \label{eq:linear-weak-transfer}
\end{equation}
Both spatial regulator corrections are retained.  If $\beta(t)X\le\pi$,
then $r_u=-2\beta u$ and $\Delta r=-2\beta$ are bounded on the full
buffer, so this is $O(e)$ even in the unresolved regulator layer.

For $\beta(t)X>\pi$ the weak-test factor need not be bounded.
Compare each actual density separately with its own zero measure
$\nu_F$.  Set
\begin{equation}
 \lambda_X(t)=\frac{\sup_{|u-X|\le\ell+2}\epsilon_g(t,u)}{A_X}.
 \label{eq:linear-regulator-ratio}
\end{equation}
The large real saddle of the protocol satisfies, for $t\ge64$,
\begin{equation}
 \tfrac14<L(64)\le L(t)\le C\log(2+t),\qquad
 ct\le\kappa(t)\le Ct\log(2+t),\qquad
 |\log E_0(t)|\le Ct\log^2(2+t).
 \label{eq:linear-regulator-growth}
\end{equation}
The saddle residual at $u=1/4$ is negative at $64$ since
$4\pi e<41$, and $L'=2L/\kappa>0$.  Taking logarithms of
$e^{4L}\le Ct(1+L)$ gives the upper bound for $L$; then
$\kappa=8tL-2t+36$ gives its bounds.  Substitute
$\pi e^{4L}=tL/2+9/4$ into $M_0,E_0$ for the last estimate.
In particular $\beta$ is globally bounded on $t\ge32$, and
$\beta\ge c/[t\log(2+t)]$ for $t\ge128$.

For $64<t<128$, $E_0$ is bounded and
\[
 \log\lambda_X\le-\tfrac14\beta X^2+\pi X/8+C\log^2X
 \le-\pi X/8+C\log^2X.
\]
For $t\ge128$ on the joint range,
$\beta\ge c/[\log X\log(2+\log X)]$, whereas
$|\log E_0|\le C\log X\log^2(2+\log X)$.
The negative quadratic regulator exponent dominates these terms and
the carrier's $-\pi X/8+O(\log X)+t\Re(\alpha^2)/4$ logarithm,
whose last term is nonnegative for large $X$.  Consequently, uniformly
in this second regime,
$\lambda_X\le Ce^{-cX}$, $|\log\lambda_X|\le CX^2$ and
$b_{r,1}+b_{r,2}\le C(1+X)$.

The physical exterior and interior floors of
Lemma~\ref{lem:linear-collars} give, for $F=H,P_{\rm lin}$,
\begin{align}
 \int_\Omega\log(1+q/|F|^2)
 &\le C\lambda_X^2\left[
 1+\frac{1+\log(1+C/\lambda_X^2)}\omega\right],\notag\\
 |R_F-K*\nu_F|
 &\le C(1+X^2)\lambda_X^2\left[
 1+\frac{1+\log(1+C/\lambda_X^2)}\omega\right]
 \le Ce^{-cX}.
 \label{eq:linear-regulator-bias}
\end{align}
Indeed each full collar uses the elementary integral
$\int_{|z|\le R}\log(1+a^2/|z|^2)\,dA
\le\pi a^2[1+\log(1+R^2/a^2)]$ with
$a\asymp\lambda_X/\omega$, and there are $O(\omega)$ collars.
The weak zero-measure identity has test allowance
\[
 \|\Delta K\|_\infty
 +2b_{r,1}\|\nabla K\|_\infty
 +(b_{r,2}+b_{r,1}^2)\|K\|_\infty,
\]
which explains the retained polynomial factor.  Root displacement
$O(e/\omega)$ and $O(\omega)$ roots then give
$|K*\nu_H-K*\nu_{P_{\rm lin}}|\le Ce$.  The two regimes prove
the density assertion; the same argument pays the compact-time rate.

For the other observables apply \eqref{eq:linear-L1-transfer}
directly.  Every weak derivative of $K$ through total order four is
bounded: the kernel is $C^3$ and its fourth derivative is bounded and
piecewise polynomial, without a boundary delta.  Transfer each
spatial derivative to $K$ and multiply its supremum by the $L^1$
error.  Since $q>0$, these weak filtered derivatives represent the
ordinary derivatives.  Evenness and conjugation of $H$, the regulator
and the kernels supply the negative-height statement.
\end{proof}
This static comparison controls neither potential time derivatives
nor the products of unfiltered gradients in the actual covariance.
It uses no regulator floor on an unbounded height tail.

\subsection{A finite phase formula in the resolved regime}
\begin{corollary}[Outer finite phase]\label{cor:outer-finite-phase}
Let $I=[t_a,t_b]\Subset(64,\infty)$.  Define
\begin{equation}
 p_{\rm out}=\frac\pi\omega,\qquad
 v_{\rm out}=\frac\theta\omega,\qquad
 R_{\rm out}(t,X,y)=k_h(y)P(p_{\rm out},v_{\rm out}),
 \label{eq:outer-finite-phase}
\end{equation}
where $P(p,v)$ is the existing exact five-coefficient periodization
\eqref{eq:finiteP}.  Uniformly for $t\in I$, $|y|\le Y$,
\begin{align}
 |R_H-R_{\rm out}|&\le C_I\frac{1+\log^2X}{X},
 \label{eq:outer-finite-phase-error}\\
 \left|R_H-\frac{k_h(y)}{4\pi}\log\frac X{4\pi}\right|
 &\le C_I\left\{(\log X)^{-4}+\frac{1+\log^2X}{X}\right\}.
 \label{eq:outer-continuum-error}
\end{align}
On $128\le t\le\tfrac14\log X$ the finite phase model in
\eqref{eq:outer-finite-phase} has the joint error
$CX^{-3/4}(1+\log^4X)$.
\end{corollary}
\begin{proof}
On fixed $I$, $\beta_*:=\min_I\beta>0$ and
\begin{equation}
 \log\lambda_X\le-\tfrac12\beta_*(X-\ell-2)^2
       +\pi X/8+C_I(1+\log^2X),\qquad
 |\log\lambda_X|\le C_I(1+X^2).
 \label{eq:linear-compact-regulator}
\end{equation}
The physical-collar payment therefore gives the stronger compact bound
\begin{equation}
 |R_H-K*\nu_H|\le C_Ie^{-\beta_*X^2/4}.
 \label{eq:linear-compact-bias}
\end{equation}
Freeze the phase roots at $X$:
\[
 b_j=X+\frac{(j+1/2)\pi-\theta}{\omega},\qquad
 |a_j-b_j|\le\frac C\omega\{\delta(X)+X^{-1}\}.
\]
The $O(\omega)$ root count and global Lipschitz continuity of the
zero-extended $k_\ell$ give
\[
 \left|K*\nu_H-k_h(y)
  \sum_{j\in\mathbb Z}k_\ell(v_{\rm out}-(j+1/2)p_{\rm out})\right|
 \le C_I\delta(X).
\]
Kernel-edge crossings cost no missing labels.  The sum is exactly
$P(p_{\rm out},v_{\rm out})$ by \eqref{eq:latticeidentity}, so its
finite formula has zero truncation error.  Bounded periodic Bernoulli
polynomials in \eqref{eq:finiteP} give
$|P(p,v)-1/p|\le C_\ell p^4$ for small $p$; this is the sharper
degree-eight-kernel rate.  Together with
\eqref{eq:linear-phase-bounds} it proves
\eqref{eq:outer-continuum-error}.  For $t\ge128$ the joint range is
uniformly in the resolved second regulator regime after one onset;
the same argument uses \eqref{eq:linear-regulator-bias} and
$e=C\mathcal E(t,X)$, while phase freezing costs $O(X^{-1})$.
\end{proof}
The finite zero-lattice formula is not applied throughout the
unresolved regulator layer.  These estimates apply to both tails by
$X=|x|$.  For $|y|>h$ the leading compact-time profiles vanish, and
the regularization payment still controls the off-axis density.
Widths and transverse strips remain fixed.

\subsection{Two regularized layers and their exact switching center}
For the linear-model amplitude and frequency in
\eqref{eq:linear-frequency}, set
\begin{equation}
 b(t,X)=\log\frac{\epsilon_g(t,X)}{A_X(t)},\qquad
 d(t,X)=\frac{b(t,X)_+}{\omega(t,X)},\qquad
 R_{\rm layer}=\frac\omega{2\pi}
       \{k_h(y-d)+k_h(y+d)\}.
 \label{eq:linear-layer}
\end{equation}
The off-axis layers are regularized density, not nonreal zeros.
When $d=0$ they merge into $\omega k_h(y)/\pi$; when $d$ is large
they leave the fixed observation strip.

\begin{theorem}[Coarse regulator-layer profile]\label{thm:regulator-layer}
Uniformly on $32\le t\le\tfrac14\log X$, $|y|\le Y$,
\begin{equation}
 |R_H(t,X,y)-R_{\rm layer}(t,X,y)|\le C.
 \label{eq:linear-layer-error}
\end{equation}
For each fixed compact $[32,T]$ the same conclusion holds with its own
$C_T$.  Its absolute error is $O(1)$, distinct from the vanishing
static error in Theorem~\ref{thm:linear-static}.  In particular
\begin{equation}
 \frac{R_H}{\log X}=
 \frac{k_h(y-d)+k_h(y+d)}{8\pi}+O(1/\log X)
 \label{eq:linear-layer-normalized}
\end{equation}
on either declared range.
\end{theorem}
\begin{proof}
When $\beta X\le\pi$, the logarithmic regulator/carrier ratio varies
by $O(1)$ on the full fixed buffer, and
$\Im\Theta_{\rm lin}(t,u+iv)=\omega v+O(X^{-1})$.
Outside fixed phase collars $|\cos\Theta_{\rm lin}|$ is comparable
to $e^{|\Im\Theta_{\rm lin}|}$.  Inside a collar the only
logarithmic singularity is bounded by
$C+2|\log|\omega(z-d_j)||$, whose integral after phase scaling is
bounded; $O(\omega)$ collars each have area $O(\omega^{-2})$.
Thus, uniformly for arbitrarily positive or negative $b$,
\begin{equation}
 \|S_{P_{\rm lin}}-2\log|\mathcal A_{\rm lin}|
                   -2\max(b,\omega|v|)\|_{L^1(\Omega)}\le C.
 \label{eq:linear-layer-potential}
\end{equation}
For $b>0$ the $q$ term supplies the lower bound; for $b<0$ the
integrable cosine logarithm does so.  The upper bound follows from
$|\cos\Theta_{\rm lin}|\le e^{|\Im\Theta_{\rm lin}|}$.
The carrier logarithm is harmonic.  The two adaptive corrections in
Lemma~\ref{lem:holomorphic-weak-density} are $O(1)$ in this regime,
since $0\le w\le1$ and the spatial derivatives of $r$ are bounded.
Finally, distributionally,
\[
 \partial_v^2\max(b,\omega|v|)=
 \begin{cases}
 \omega(\delta_d+\delta_{-d}),&b>0,\\
 2\omega\delta_0,&b\le0.
 \end{cases}
\]
Test against $K$ and use its unit horizontal integral, retaining the
horizontal average exactly.  This proves the layer formula for
$P_{\rm lin}$, and Theorem~\ref{thm:linear-static} transfers it to $H$.
For $\beta X>\pi$, the preceding proof gives $b<0$ eventually;
zero-measure transfer and finite phase yield the stronger estimate
\[
 R_H=\omega k_h(y)/\pi
       +O(\mathcal E+X^{-1}+\omega^{-4}+e^{-cX}).
\]
Both regimes give \eqref{eq:linear-layer-error}; then
$\omega/\log X=1/4+O(1/\log X)$ proves
\eqref{eq:linear-layer-normalized}.
\end{proof}

\begin{corollary}[Center and resolved switch scaling]\label{cor:switch-scaling}
For sufficiently large $X$ there is a unique $t_c(X)$ in a fixed small
right neighborhood of $64$ satisfying the exact known-model equation
\begin{equation}
 b(t_c(X),X)=0,\qquad
 \epsilon_g(t_c(X),X)
   =|\mathcal G_{t_c(X)}((1+iX)/2)|.
 \label{eq:linear-exact-center}
\end{equation}
With $\kappa_0=\kappa(64)$,
\begin{equation}
 t_c-64=\frac{64}{\log X+c_0+O(1/\log X)},\qquad
 c_0=\log\frac{4e}{\pi\kappa_0}.
 \label{eq:linear-center-location}
\end{equation}
For bounded $\tau$, retain this exact $t_c$ and put
\begin{equation}
 t_X(\tau)=t_c(X)-\frac{128\tau}{\pi X\log X},\qquad
 \tau_+=\max(\tau,0).
 \label{eq:linear-switch-time}
\end{equation}
Uniformly for bounded $\tau$ and $|y|\le Y$,
\begin{equation}
 \frac{R_H(t_X(\tau),X,y)}{\log X}
 =\frac{k_h(y-\tau_+)+k_h(y+\tau_+)}{8\pi}
       +O(1/\log X).
 \label{eq:linear-switch-law}
\end{equation}
\end{corollary}
\begin{proof}
Bounded-real-part Stirling near $t=64$ gives
\begin{equation}
 \log A_X(t)=-\pi X/8+\tfrac74\log X
       +\frac t{16}\log^2(X/(4\pi))+O(1).
 \label{eq:linear-center-amplitude}
\end{equation}
Write $\Delta=t-64>0$.  Direct differentiation of the actual flat
switch, including the denominator in $\chi$, and smoothness of the
saddle give
\begin{equation}
 \log\beta(64+\Delta)
   =-64/\Delta+1-\log\kappa_0+O(\Delta),\qquad
 \frac{\beta'}\beta=64/\Delta^2+O(1).
 \label{eq:linear-switch-derivatives}
\end{equation}
The directly computed time derivative of the known ratio is
\begin{equation}
 b_t=(\log E_0)'-\tfrac12\beta'X^2
                         -\tfrac14\Re\{\alpha(s_X)^2\}.
 \label{eq:linear-ratio-time-derivative}
\end{equation}
Near $64$, $\beta'>0$, the first term is bounded and the last is
negative of order $\log^2X$.  Hence $b$ is strictly decreasing there
for large $X$.  At $64$ it is positive of order $X$; at any fixed
nearby time above $64$ it is negative of order $X^2$.  This proves the
stated local existence and uniqueness, without an all-time uniqueness
claim.  At the exact center,
\begin{equation}
 \begin{aligned}
 \beta(t_c)&=\frac\pi{4X}+O(\log^2X/X^2),&
 \omega(t_c,X)&=\tfrac14\log X+O(1),\\
 b_t(t_c,X)&=-\frac{\pi X\log^2X}{512}
                   (1+O(1/\log X)).
 \end{aligned}
 \label{eq:linear-center-derivatives}
\end{equation}
Taking logarithms of the first relation and using
\eqref{eq:linear-switch-derivatives} proves
\eqref{eq:linear-center-location}.  The fixed $\eta$ and amplitude
constants enter finer terms, not $c_0$.
On an interval of width $O(1/(X\log X))$ about $t_c$, the same direct
switch derivatives give
\begin{equation}
 b_{tt}=O(X\log^4X).
 \label{eq:linear-switch-second-derivative}
\end{equation}
This is not differentiation of the error in
\eqref{eq:linear-center-location}. Taylor's formula and
\eqref{eq:linear-center-derivatives} now yield
\[
 \frac{b(t_X(\tau),X)}{\omega(t_X(\tau),X)}
       =\tau+O(1/\log X).
\]
Insert this into Theorem~\ref{thm:regulator-layer} and use Lipschitz
continuity of $k_h$ to prove \eqref{eq:linear-switch-law}.
\end{proof}
The exact implicit center in \eqref{eq:linear-switch-time} cannot be
replaced by $64+64/\log X$ or by a short inverse-logarithm expansion:
their omitted time error is much larger than the width
$1/(X\log X)$.  Its amplitude is specifically that of
$P_{\rm lin}$; substituting a full-saddle amplitude would require a
separate recentering estimate.  No numerical value of this center is
certified here.

For each fixed $y$ with $|y|<h$, the pre-switch tail and
\eqref{eq:outer-continuum-error} also give the nonuniform fixed-time
limit
\begin{equation}
 \lim_{X\to\infty}\frac{R_H(t,X,y)}{\log X}
 =\begin{cases}
 0,&t\le64,\\
 k_h(y)/(4\pi),&t>64.
 \end{cases}
 \label{eq:linear-fixed-time-jump}
\end{equation}
Every finite-height expression is continuous in time, so convergence
to this displayed limit is not uniform on a time interval crossing
$64$.  The chosen regulator causes the nonuniform limit; $H$ has no
finite-height time discontinuity.  This does not obstruct a time/height
interpolation and gives no RH or earlier-simplicity conclusion.

\subsection{Shift, rounding and the open correlation question}
\label{sec:crossover-outlook}
The quantities controlling the existing result are, together,
\begin{equation}
 b=\log(\epsilon_g/A_X),\quad b(t_c,X)=0,\quad d=b_+/\omega,
 \qquad R_{\rm layer}=\frac{\omega}{2\pi}
          [k_h(y-d)+k_h(y+d)].
 \label{eq:crossover-interpretation}
\end{equation}
The center shift $t_c-64\sim64/\log X$ and the rounding width
$1/(X\log X)$ differ greatly. Fisher--Barber's distinction between
location shift and rounding motivates this comparison \cite{fisherbarber};
no film-model exponent or thermodynamic interpretation is imported.
Corollary~\ref{cor:switch-scaling} is uniform for bounded scaled time
$\tau$ and $|y|\le Y$ on its stated full buffers, using the exact implicit
center. Its normalized $O(1/\log X)$ remainder is an unnormalized $O(1)$
allowance, not a vanishing absolute error for the global contract.
No time/current/covariance theorem follows by differentiating that allowance.
The two regularized layers are not nonreal zeros.

One possible diagnostic is
\begin{equation}
 m_X(t)=\frac{4\pi R_H(t,X,0)}{k_h(0)\log X}.
 \label{eq:crossover-diagnostic}
\end{equation}
Its fixed-time limits are zero for $t\le64$ and one for $t>64$, by
\eqref{eq:linear-fixed-time-jump}. It does not replace the unnormalized
target or establish an intrinsic Landau order parameter. Note~38
\cite{n38} treats statistical-field-theory language as a possible method,
with missing premises, rather than an established theta structure.

For a declared observable $O$ and normalized window $W_X$, a possible
connected statistic is
\begin{align}
 \mathcal C_X(s;t)={}&\langle O(t,X+u)O(t,X+u+s)\rangle_{W_X}\notag\\
 &-\langle O(t,X+u)\rangle_{W_X}
           \langle O(t,X+u+s)\rangle_{W_X}.
 \label{eq:connected-outlook}
\end{align}
A future specification must choose $O$, centering, amplitude scales,
window width and limiting regime, then define a perturbation and response.
The signed local covariance $Q$ is not this statistic. An integrated
correlation is not automatically a susceptibility: that interpretation
needs a response identity and its hypotheses. Correlation-length growth
and integrated response are relevant tests; a fixed-separation value need
not diverge. Tong's order-parameter and connected-correlation discussion
\cite{tong} supplies methodological context, not a theta theorem.
No Gibbs structure, RG fixed point, conformal invariance or intrinsic
criticality is established. The scoped literature inspection and its
limits are recorded in the
\repo{references/PHASE_TRANSITION_METHODS_20261024.md}{primary-literature record}.

\subsection{The exact residual and the historical spatial restriction}
The linear pair is not an exact heat solution.  Differentiating its
explicit factors gives
\begin{align}
 (\partial_t+\partial_{zz})P_{\rm lin}
 &=\frac14\sum_{w=s,1-s}\mathcal G_t(w)
       \{\alpha(w)^2-D_t(w)^2-D_t'(w)\},\notag\\
 \alpha^2-D_t^2-D_t'
 &=-\alpha'-t\alpha^2\alpha'
       -\frac t2\{(\alpha')^2+\alpha\alpha''\}
       -\frac{t^2}4\alpha^2(\alpha')^2.
 \label{eq:linear-heat-residual}
\end{align}
In the proved joint range,
\begin{equation}
 |(\partial_t+\partial_{zz})P_{\rm lin}|
 \le C\frac{1+t\log^2X}{X}
        \{ |\mathcal G_t(s)|+|\mathcal G_t(1-s)|\}.
 \label{eq:linear-heat-residual-bound}
\end{equation}
This nonzero residual is a field budget, not a propagated evolution
error for current, density time derivative, forcing, potential time
derivative or the actual $Q,\Delta Q$.  The static proof uses the
holomorphic weak identity rather than imposing the exact heat-field
evolution equations on $P_{\rm lin}$.

At fixed large $t$, the sufficient joint domain above requires
$X\ge e^{4t}$; it has no proved overlap with the polynomial inner
window $X\asymp\kappa^{1/3}$ of Theorem~\ref{thm:explicit}.
That restriction comes from the absolute Gaussian integration loss;
it is not a necessary limitation of all gamma representations.
The next section removes this spatial restriction and proves static
inner overlap using a different, full-saddle field.
